\section{Constructing the comparison values over the coefficient field}
\label{sec:value-existence}

\status{VE1--VE9 remove the separate polynomial-value existence premise
from CS. They prove block rationality, adelic point existence and
representative independence, then construct the literal \texttt{ValEx}
and concrete comparison under supplied global forms, standard-frame,
level and character data. FC19 now supplies those data from the actual
true-column and initial-realization hypotheses, and VE10 obtains their
concrete comparison. DR1--DR3 then construct that initial realization
from the supplied CM and decoration data; VE11 discharges this remaining
geometric premise. The independent metric, spectral, cohomological
and arithmetic-window inputs retain their separate scope.}

\paragraph{The data and the assertion.}
Let $E/F$ be the CM extension, $L$ the coefficient field with its
specified embedding $\jmath:L\hookrightarrow\mathbb C$, and
$t_\tau:E\hookrightarrow L$ the specified lifts of the complex
embeddings. Keep the same frame $B$, place identification and position
map as LI and SF. The supplied level-zero forms give functions
$f_{x,c}(z)$ for $x\in V$, a class $c$ and an invertible complex
matrix frame $z=(z_\tau)_\tau$. Here the source's
\texttt{Frames} consists of all invertible matrix frames: its Gram
label does not require the matrices themselves to be unitary.

For an away component $i$ and a tuple of full local lattices
$\mathbf L$, we will construct a unique polynomial
$Q_{x,i,\mathbf L}$ over $L$ in the prescribed tensor representation.
For every basis $g$ of $\mathbf L$ and every presentation of its
associated adelic point as $a=\gamma y_c\kappa$, its evaluation must
equal
\[
 \operatorname{ev}_C(\jmath Q_{x,i,\mathbf L})(z)
       =\chi(\kappa)f_{x,c}(z'),\qquad z=\gamma_\infty z'.
                                                        \tag{VE.1}
\]
Thus one polynomial must work for \emph{all} these presentations.

\paragraph{VE1: explicit rational polynomials for each block.}
Write $k=m_{\rm loc}+1$, $\lambda=(\lambda_1,\ldots,\lambda_k)$,
$e=\lambda_k$, $\lambda^*=(-\lambda_k,\ldots,-\lambda_1)$ and
$e^*=-\lambda_1$. Put
\[
 T(Y)=Y^{-\mathsf T}R_k,
\]
where $R_k$ reverses columns. The integral duality of CS2 supplies,
for every $f\in\Delta_{\nu(\lambda)}(\mathbb Z)$, a polynomial
$f^*\in\Delta_{\nu(\lambda^*)}(\mathbb Z)$ satisfying
\[
 \det(Y)^e f(Y)=\det(T(Y))^{e^*}f^*(T(Y))
       \quad(Y\in\mathrm{GL}_k).
                                                        \tag{VE.2}
\]
Its sign normalization is exactly the one in CS2; explicitly,
$f^*(X)=\epsilon^{\lambda_k}\det(X)^{\lambda_1-\lambda_k}
f(X^{-\mathsf T}R_k)$ on invertible matrices, with
$\epsilon=\det R_k$. Complementary minors show that this is an
integer polynomial on the specified Weyl space.

Suppose the block coordinate matrix is $M=\jmath(M_0)$ with
$M_0\in\mathrm{GL}_k(L)$. Define
\[
 F_{\rm ev}(X)=\det(M_0)^{-e}f(M_0^{-1}X),\qquad
 F_{\rm odd}(X)=\det(M_0)^{e^*}f^*(M_0^{\mathsf T}X).
                                                        \tag{VE.3}
\]
All coefficients lie in $L$. These polynomials belong to the
required dual-Weyl block spaces. Indeed integral Weyl generators lie
in the column-homogeneous, upper-column-unipotent invariant space;
row substitutions commute with the column action and preserve
column degrees. Scalar multiplication preserves this space too.
No equality of the integral lattices $\Delta$ and $\nabla$ is used.

If $Z=MY$, the even expression gives
$\det Z^e\,\jmath(F_{\rm ev})(Z)=\det Y^e f(Y)$.
If $Z=T(MY)=M^{-\mathsf T}T(Y)$, then
$M^{\mathsf T}Z=T(Y)$, and (VE.2) gives
\[
 \det Z^{e^*}\,\jmath(F_{\rm odd})(Z)=\det Y^e f(Y).
                                                        \tag{VE.4}
\]
The source's \texttt{WtTie} assigns precisely these two weights to
even and odd positions. Equations (VE.3)--(VE.4) therefore prove
\texttt{BlockRationalityField}. For the actual frame coordinates,
$M_0=t_\tau(B^{-1})$; hence they also prove
\texttt{BlockRationalityAtPlace}, with the original embeddings.

\paragraph{VE2: tensor the blocks and recover representative values.}
The rationality field of the supplied \texttt{Forms0} says that each
$f_{x,c}$ is a finite sum of terms
\[
 \jmath(r)\prod_\tau\det(z_\tau)^{\lambda_{\tau,k}}
                                  f_\tau(z_\tau),
 \qquad r\in L,\quad f_\tau\in\Delta_{\nu(\lambda_\tau)}(\mathbb Z).
                                                        \tag{VE.5}
\]
For one term choose the block polynomials from VE1, multiply them
in disjoint variable blocks, and multiply by $r$. The result belongs
to the source's $L$-tensor space \texttt{VtupR}. Evaluating its product
gives (VE.5), since the position map is a bijection. Add the resulting
polynomials for a finite sum, and use their negatives for additive
inverses. This proves \texttt{evRat\_gen}, \texttt{evRat\_of\_blocks}
and the actual \texttt{EvRatRep}: every representative value has a
polynomial $Q_c$ over $L$.

\paragraph{VE3: rational transport from one class representative.}
For a rational unitary $\gamma$, put $H=B^{-1}\gamma B$.
The same standard-frame identity as SF gives $H^*H=I$. In position
$(v,s)$ define the coefficient-field matrix
\[
 h_{v,s}=
 \begin{cases}
 t_\tau(H),&s\text{ even},\\
 t_\tau(\overline H),&s\text{ odd},
 \end{cases}
 \qquad\tau=\operatorname{pos}(v,s).
                                                        \tag{VE.6}
\]
Because $\overline H=(H^{-1})^{\mathsf T}$, these are the exact
even and odd frame-transport matrices, including the conjugation
already required by LI. The natural action is
\[
 (\rho(h)Q)(X)=
    \left(\prod_{v,s}\det(h_{v,s})^{-e_{v,s}}\right)
             Q\bigl((h_{v,s}^{-1}X_{v,s})_{v,s}\bigr).
                                                        \tag{VE.7}
\]
It preserves \texttt{VtupR}: each row substitution preserves its
dual-Weyl block, tensor generators map to tensor generators, and
linearity handles their span. This proves the source's
\texttt{VtupStab} directly.

If $z=\gamma_\infty z'$, then the block coordinates satisfy
$Z(z)=hZ(z')$. Substitution in (VE.7), including its determinant
factor, proves
$\operatorname{ev}_C(\jmath\rho(h)Q)(z)
=\operatorname{ev}_C(\jmath Q)(z')$.
The supplied level-zero character rationality gives
$\chi(\kappa)=\jmath(r_\kappa)$: the level character is a finite
product of the supplied character values and their inverses, and
the higher-block factors are absent at level zero. Thus
\[
 Q=r_\kappa\rho(h)Q_c
                                                        \tag{VE.8}
\]
is a rational polynomial giving the right side of (VE.1) for this
presentation. This proves \texttt{EvRatPt}; the rational action and
character lift have been constructed rather than postulated.

\paragraph{VE4: construct the required adelic points.}
At each place above $p$, LI's split unitary group is isomorphic to
$\mathrm{GL}_k(F_v)$ by projection to its first component; the
second component is determined by the Hermitian equation. Given
$g=(g_v)_v$, insert the first components $B_vg_vB_v^{-1}$ at these
finitely many places and the identity at all others. The resulting
section $j_p(g)$ is a finite adele in the unitary group.
The away representative $y_i$ is the identity above $p$, so
\[
 a=y_i j_p(g)
                                                        \tag{VE.9}
\]
has the exact prescribed local components and equals $y_i$ away
from $p$. The supplied class representatives exhaust the double
quotient; hence $a=\gamma y_c\kappa$ for a rational unitary
$\gamma$ and $\kappa\in U$. This is \texttt{PtEx}.

The same finite-place section also proves the two companion point
properties used by CS. For any unitary adele $b$, divide it by the
section of its components above $p$. The quotient is the identity
above $p$ and equals $b$ elsewhere. If $b$ is in the level away
from $p$, this quotient lies in the full level. Applying this to
the away double-coset decomposition of $y_c$ gives \texttt{PtAll}.
Applying it to $y_i^{-1}\gamma y_i$ for an away stabilizer gives
\texttt{AwayEx}. This uses finite adelic insertion and the supplied
class decompositions; it does not require a new strong-approximation
theorem for rational unitary points.

\paragraph{VE5: the ambiguity is in the kernel of the character.}
Suppose $a,a'$ describe the same $i,\mathbf L$, with bases $g,g'$.
Then $\kappa_0=a^{-1}a'$ is the identity away from $p$. Above $p$,
its first split component in the common frame is $g_v^{-1}g'_v$.
Equality of the lattice spans says that this matrix and its inverse
are integral. The actual hyperspecial description of the level,
proved in LI/SF and CR3, puts $\kappa_0$ in every local level group,
hence in $U$. The good-prime condition excludes the two auxiliary
character primes. Since $\kappa_0$ is the identity there, the
literal level-character support condition gives
\[
 a'=a\kappa_0,\qquad \chi(\kappa_0)=1.
                                                        \tag{VE.10}
\]
This proves \texttt{PtOne}. The equality of local level groups here
is \texttt{Joint0.Uloc\_eq}; it is not inferred from rank equality.

\paragraph{VE6: one polynomial works for every presentation.}
The supplied \texttt{Forms0} also gives an adelic extension
$F_x(a,z)$ of $f_{x,c}(z)$, with
\[
 F_x(\gamma y_c\kappa,z)
       =\chi(\kappa)f_{x,c}(z')\quad(z=\gamma_\infty z').
                                                        \tag{VE.11}
\]
This is \texttt{AdExt}: membership in the actual form space and its
adelic realization are fields of the supplied global form datum.
Choose one basis and one presentation by VE4 and its polynomial
by VE3. Its evaluation is $F_x(a,z)$ for every $z$.
For another basis and presentation, VE5 gives $a'=a\kappa_0$ with
$\chi(\kappa_0)=1$, so $F_x(a',z)=F_x(a,z)$. Applying (VE.11)
to the second presentation proves (VE.1) for the same polynomial.
This proves \texttt{valEx\_of\_rep'} and the actual
\texttt{ex\_C\_blk}. All basis and representative choices disappear
from the required evaluation identity.

\paragraph{VE7: uniqueness and linearity require no extra existence.}
The block-coordinate map from frames onto tuples of invertible
matrices is surjective: the even inverse is $Y=M^{-1}Z$, and the
odd inverse is
$Y=M^{-1}(ZR_k)^{-\mathsf T}$.
The determinant factors in evaluation are nonzero. Consequently a
polynomial with zero frame evaluations vanishes on all invertible
block tuples. To see that it is zero everywhere, let $\mathfrak D$
be the product of their generic determinants. The polynomial
$Q\mathfrak D$ vanishes on every tuple: either every block is
invertible or one factor of $\mathfrak D$ is zero. Polynomial
functions over the infinite field $\mathbb C$ separate polynomials,
and $\mathfrak D\ne0$ by evaluation at the identity tuple.
Thus $Q=0$.

If two polynomials satisfy (VE.1), choose one presentation by VE4
and subtract their evaluations. The preceding argument and the
injectivity of $\jmath$ prove equality over $L$. The unique value
$Q_{x,i,\mathbf L}$ is additive and $L$-linear in $x$, since sums
and scalar multiples satisfy the same identity. After mapping
coefficients by $e:L\hookrightarrow K$, and extending by zero on
non-lattice tuples as in the source, this constructs the literal
semilinear comparison \texttt{cmpL}.

\paragraph{VE8: supply all fields of the concrete comparison.}
For the actual recursion, the input data are
\texttt{Forms0}, the common standard $B$, its matching local frames,
\texttt{Joint0}, the coefficient-field level character, and LI's
original position/valuation compatibility. VE1--VE2 supply block
and representative rationality; VE4--VE5 give the point properties;
VE6--VE7 give existence, uniqueness and evaluation injectivity.
LI8 supplies the matching inverse away character. CS1--CS4 now apply
with a \emph{proved} \texttt{ValEx}, so supply \texttt{ValIntRec};
CR and CS5--CS9 give the exact integral span. Hence every field of
\texttt{CmpHyp0P} is supplied:
\[
 \texttt{ex},\ \texttt{evInj},\ \texttt{ptEx},\
 \texttt{ptAll},\ \texttt{chi},\ \texttt{latM}.
                                                        \tag{VE.12}
\]
In particular this reconstructs \texttt{comparison\_statements}
and \texttt{exists\_isConcrete}. The temporary use of CS under a
\texttt{ValEx} premise is now discharged without circularity: VE1--VE7
do not use CS's span or integrality conclusions, only its independent
integral duality CS2.

\paragraph{VE9: the exact conditional existence theorem.}
Suppose a supplied $D7$ has level data $L_0$, forms $T_u$, a common
standard-frame datum, \texttt{Joint0} for each $u$, coefficient-field
level-zero characters, and companion ideals
\[
 L_{b,u}+\mathfrak a_u L_{b,u}^{\#}=\Lambda_0.
                                                        \tag{VE.13}
\]
For every supplied scale $\mathrm{sc}$, construct the local
identifications and their original coefficient embeddings by LI;
construct the inverse away characters by LI8; replace only the frames
using SF and (VE.13); and supply the ranks by CR. The resulting
record uses $C_e=1$, exactly as in
\texttt{exists\_concrete\_of\_record}. Its forms are transported
unchanged to the recursion, so the joint level and character
identities are retained. VE8 then gives a concrete comparison $c$.
For every good place its actual image satisfies
\[
 \boxed{\quad\mathcal O\text{-span}\,c(\mathcal L)=M.\quad}
                                                        \tag{VE.14}
\]
This closes the former extra block-rationality and polynomial-value
premises. The remaining construction problem is to supply the
global forms, level/character and standard companion data from
the same true-column geometry. The chart and form constructions
in FC and CV address that problem. None of
(VE.12)--(VE.14) constructs spectral estimates, metric transport,
cohomological windows, or a global prime-sum estimate.

\paragraph{VE10: the actual true-column comparison follows from the direct form construction.}
Take the actual true-column record $I:\texttt{TC}(B,X)$ with its
specified witnesses $I.\texttt{Ok}$, $\texttt{Pins}$,
$l:I.\texttt{Laws}$ and $Z:\texttt{Later}(B,I.\texttt{cols}(l))$.
Put $D=I.\texttt{rec7}\,l\,Z$. Retain the explicit geometric
hypothesis that the valid framed decorated points used in the chart
and neighbor constructions have unitary finite-adelic realizers.
Then for every supplied scale the actual concrete comparison exists,
with $C_e=1$ and with its original coefficient embeddings and positions.
More precisely, the output is the recursion
$\operatorname{recC}(D,\mathrm{sc},1,L_0,I_{\rm loc},\psi,h_d)$
constructed by \texttt{exists\_concrete\_of\_record}, with its transported
forms and concrete comparison; it is not a change to the definition of $D$.

Indeed FC19 constructs $T_u:\texttt{Forms0}\,D\,u$ for every
word, using the actual CV values, NB up/down naturality and FC residue,
norm, distinguished-block and dilation fields. It simultaneously
constructs $L_0(u)$ with the exact five \texttt{Joint0} carrier
equalities. Its \texttt{StdTied} data use the same standard coordinates;
FC16 gives the coefficient-field level-zero character values; and
the companion ideal is the explicit $(\bar\pi_u)$, with the exact
identity (FC.41). These are every global-record premise of VE9.
Apply VE9 to these objects, rather than to a separately supplied
form record. At each good place $w$ the resulting map therefore obeys
\[
 \mathcal O_w\text{-span}\,
          c_{u,w}(\mathcal L_u)=M_{u,w},\qquad C_e=1. \tag{VE.15}
\]
Here $\mathcal L_u$ is the actual global forms lattice and $M_{u,w}$
is the actual target coefficient lattice of the local model. The
assertion is an integral span equality; no index is suppressed.
Existence, uniqueness, integrality and the rank identification are
the already proved VE and CS/CR statements on these same objects.

In this intermediate formulation the initial realizer is an explicit
geometric supplier, constructed in DR and applied in VE11 below. AD1
constructs its diagonal base point, AD4 supplies global-lattice
recovery and exact frame correction, and AD3 supplies its supported
modifications; none of these manufactures the initial decorated
realizer. The full terminal \texttt{FormsTiedSym}/\texttt{ColsTiedSym}
construction is a separate remaining task and is not used in this
direct level-zero comparison. The true-column arithmetic laws and
the later metric, spectral, cohomological and prime-window arguments
retain their stated boundaries.

\paragraph{VE11: discharge the initial-realizer premise on the actual column data.}
Keep exactly the true-column data and scale in VE10. Its CM-field
record includes the literal unramified condition (DR.1), and its
valid framed decorations have the self-duality, nonzero type
functionals, auxiliary residue bases and rational local frames used
in DR3. The type, auxiliary and twist prime sets satisfy (DR.12):
the type primes divide the selected $T_u$, the others do not;
the split-level data orient the auxiliary/twist union; and
$(\pi_u,\bar\pi_u)=\mathcal O$ orients the type primes.
The disjointness checks are the ones in
\texttt{exists\_ptOf\_twists}. Thus DR3 supplies a unitary finite-adelic
realizer for every valid framed decorated point occurring in the
charts and neighbor moves.

There is no stronger rational-genus hypothesis at this step.
DR2 proves the source's actual adelic \texttt{Col.UnimodGenus}
by constructing the matrix and its carried lattice. DR3 installs
all the residue and frame clauses of \texttt{IsPtOf}. The rational
semilocal predicates with similar genus names are bypassed by this
construction. In particular no rational unitary approximation or
global integral-basis assumption is inserted.

Apply VE10 with these supplied realizers. For every supplied scale
$\mathrm{sc}$ there exist the required local identifications, levels,
characters and rank data for
\[
 \operatorname{recC}(I.\texttt{rec7}\,l\,Z,
                  \mathrm{sc},1,L_0,I_{\rm loc},\psi,h_d),
                                                        \tag{VE.16}
\]
and its concrete comparison satisfies (VE.15) at every good place.
The form spaces, coefficient embeddings, odd-position involutions
and normalized valuations are precisely the ones retained in VE10.
This removes the independent decorated-point-realization premise
from the actual integral comparison theorem.

The dependency is one-way: DR1 uses local unramified Hermitian algebra;
DR2 uses only local duality, lattice recovery and finite adelic assembly;
DR3 uses exact frame modification. None uses the value polynomials,
comparison maps or their integral-span conclusion. The supplied
arithmetic host and actual $I.\texttt{Ok}/\texttt{Pins}/
\texttt{Laws}/\texttt{Later}$ remain hypotheses. The full terminal
symmetric-column, metric, spectral, cohomological and prime-window
constructions remain the separate tasks already identified.
